\section{Related Work}

\subsection{World Action Models}
World-Action Models (WAMs) are models that combine world modelling and action prediction into a tightly coupled architecture. Some previous works have described their approach as using Video-Action Models~\cite{MimicVideo, UVAM}, Vision-Language-Action models with a world modelling objective\cite{VLAJEPA, won2025dual}, or Video-Language-Action models~\cite{GR-2}, however the category of World-Action models, as defined in \citet{DreamZero}, encapsulates all these approaches. The common point that ties them together is their utilization of a future state predicting objective alongside action prediction. Although the most common to date, video is just one world modelling modality; future work, including our proposed method, may utilize alternate predictive modalities, such as learned latent represenations, 3D pointclouds, tactile senses, auditory signals, or a multimodal combination of others.

World-Action Models present a distinct strength when compared to the more prevalent Vision-Language-Action Models (of the non-world modelling variety); the ability to learn from diverse data sources. Most VLA architectures are constrained by the prohibitive cost of collecting successful expert demonstrations in a large variety of settings. WAMs on the other hand, are more natively able to learn from diverse environmental interaction data, and have been shown to improve multi-task performance, sample efficiency, and generalization to novel scenes and objects~\cite{DreamZero, MimicVideo, UWM, TRIWAM, CosmosPolicy}.

These strengths come at a cost, since many current World-Action models rely on large-scale video diffusion models for their dynamics priors, often being prohibitively expensive to train and deploy in real time or on edge hardware~\cite{DreamZero}. 

%Although our proposed method forgoes these video generation priors, we build on the strong latent representations of V-JEPA2 to reduce training and inference cost, while aiming to retain the sample efficiency and generalization benefits characteristic of WAMs.

\subsection{Vision Language Action Models}
Foundational to the recent successes in robot learning, Vision Language Action models have paved the way for generalizable and robust systems that can learn dozens or even hundreds of behaviours~\cite{pi0.5, GrootN1, geminirobotics, Octo} within a single policy. VLAs have shown great promise, but suffer from the previously mentioned limitations of data inefficiency, requiring hundreds to thousands of hours of largely repeated demonstrations to effectively generalize\cite{pi0.5, geminirobotics}. Co-training VLAs on non-robot web data improves performance in some cases~\cite{pi0.5}, but its effect is not nearly as pronouced as in WAMs.

Many leading VLAs deal with the same high computational costs as WAMs. This has motivated works like \citet{SmolVLA} and \citet{TinyVLA} to design compact VLA architectures, allowing for more accessible training and deployment, while maintaining competetive performance with leading models. To the best of our knowledge, such small scale approaches have yet to be replicated in WAMs, motivating our own approach. 

% https://arxiv.org/pdf/2507.05331

\subsection{Joint Embedding Prediction (JEPA)}
Joint Embedding Predictive Architectures, or JEPAs, are a series of world modelling architectures that prioritize semantic understanding over visual fidelity, improving efficiency over generative approaches, while acheiving comparable or superior performance on downstream tasks~\cite{vjepa2, vjepa2_1, IJEPA}. Characteristically, they differ from generative approaches like VAEs or diffusion models by exclusively operating in latent space. 

Recent works using video JEPA models have shown promising results in robotic control tasks~\cite{vjepa2, vjepa2_1, LeWorldModel}. Rather than an action-generative approach ($p(z_{i+1}, a_{i+1}| z_i)$), most JEPA works have adopted an action-conditioned formulation ($p(z_{i+1} | z_i, a_{i+1})$). This approach does have some distinct strengths (e.g. controllability), but it requires expensive inference-time planning and lacks language conditioning, limiting its practical applicability. Our proposed method instead aims to inherit the semantic and dynamic understanding properties of video JEPA models while adopting the more common action generative approach used throughout VLAs and WAMs.

\subsection{Latent World-Action Models}

%Many video and image generation models have adopted a latent-space diffusion approach to improve efficiency over pixel space diffusion~\cite{wan, latentdiffusion, CosmosPredict2.5}. Consequentially, many WAMs that build off of these generative models also adopt a latent diffusion formulation~\cite{DreamZero, MimicVideo}. However, there is a fundamental mismatch between the requirements of video generation models and WAMs, namely visual fidelity.

Although improved video generation quality has been shown to directly improve action prediction performance~\cite{DreamZero, MimicVideo}, we argue that the reliance on latents produced by Video VAEs introduces unnessecary bulk into WAMs, since they encode fine grained visual details (eg. lighting, fine grained textures) that may not be relevant to robotic manipulation tasks.

Similarly, VLA-JEPA~\cite{VLAJEPA} aims to mitigate this reliance on pixel-space VAEs by adopting a V-JEPA2 encoder alongside a generative world-action modelling objective, with promising results. However, the work mainly focuses on latent-action pretraining, and seperates their autoregressive transformer latent world model and flow matching action head. Furthermore, VLA-JEPA adopts a signficantly larger VLM backbone, V-JEPA2 encoder, world model, and action head than we do. Consequentially, their model utilizes the same amount, if not more parameters than leading VLAs\cite{pi0.5}. [TODO: calculate exact parameter count of VLA-JEPA since they dont share it]

In contrast to \citet{VLAJEPA}, our proposed method does not focus on latent-action pretraining, and unifies the world model and action head into a single DiT. We also target a much smaller scale regime, similar to \citet{SmolVLA} and \citet{TinyVLA}
